\documentclass[11pt]{article}
\usepackage[T1]{fontenc}
\usepackage{lmodern}
\usepackage[margin=1in]{geometry}
\usepackage{amsmath,amssymb,amsthm,mathtools,booktabs,microtype}
\usepackage[colorlinks=true,linkcolor=blue,citecolor=blue,urlcolor=blue]{hyperref}
\hypersetup{pdftitle={Two shorter intervals of even Dirichlet beta values},
 pdfauthor={twinheterodox}}
\newtheorem{theorem}{Theorem}
\newtheorem{proposition}[theorem]{Proposition}
\newtheorem{lemma}[theorem]{Lemma}
\theoremstyle{remark}
\newtheorem{remark}[theorem]{Remark}
\newcommand{\Z}{\mathbb Z}
\newcommand{\Q}{\mathbb Q}
\newcommand{\ind}{\mathbf1}
\newcommand{\vp}{v_p}
\DeclareMathOperator{\lcm}{lcm}
\numberwithin{equation}{section}
\title{\fontsize{16}{20}\selectfont
 \mbox{Two shorter intervals of even Dirichlet beta values}}
\author{twinheterodox}
\date{17 November 2026}

\begin{document}
\maketitle
\begin{abstract}
We prove that each of the sets
$\{\beta(6),\beta(8),\ldots,\beta(66)\}$ and
$\{\beta(8),\beta(10),\ldots,\beta(98)\}$ contains an irrational number.
These sets contain $31$ and $46$ values, respectively, improving the
previous intervals with the same lower endpoints, which contained
$35$ and $52$ values. The two conclusions concern overlapping sets;
they do not assert that the irrational values are distinct.
\end{abstract}

\section{The results and their source}

For real $s>0$, the Dirichlet beta function is
\[
 \beta(s)=\sum_{k=0}^{\infty}\frac{(-1)^k}{(2k+1)^s}.
\]
An irrationality statement about a finite set of its even values
locates an irrational value without identifying which member it is.
The lower endpoint matters: a result for a set containing $\beta(4)$
does not locate an irrational value beyond $\beta(4)$.

\begin{theorem}\label{thm:main}
Each of the following sets contains an irrational number:
\[
 \{\beta(6),\beta(8),\ldots,\beta(66)\},\qquad
 \{\beta(8),\beta(10),\ldots,\beta(98)\}.
\]
\end{theorem}

The powered-kernel construction of \cite[Theorem~1]{powered} gives
the corresponding upper endpoints $74$ and $110$. The comparison is
\begin{center}
\begin{tabular}{@{}lcc@{}}
\toprule
Lower endpoint & Previous interval & Present interval\\
\midrule
$\beta(6)$ & $6,8,\ldots,74$: $35$ values
           & $6,8,\ldots,66$: $31$ values\\
$\beta(8)$ & $8,10,\ldots,110$: $52$ values
           & $8,10,\ldots,98$: $46$ values\\
\bottomrule
\end{tabular}
\end{center}
Here the displayed integers are the arguments of $\beta$.
Neither conclusion follows just by deleting the first value from
the fifteen-value result for $\beta(4),\ldots,\beta(32)$ in
\cite{fifteen}. Similarly, the first conclusion of
Theorem~\ref{thm:main} does not imply the second: its irrational
member could be $\beta(6)$.

The proof combines two existing constructions. We use the
order-$31$ rational kernel and its isolated-pole denominator estimate
from \cite[Sections~2--4 and Appendix~A]{fifteen}, in place of the
order-$35$ kernel in \cite{powered}. We then apply the powering and
differentiation procedure of \cite[Sections~2--4]{powered}.
Squaring the kernel and differentiating four times leaves the even
indices $6,\ldots,66$; cubing it and differentiating six times leaves
$8,\ldots,98$.

Two points must be checked for this combination. First, the
denominator saving at an isolated pole must remain available after
powering. Second, the finite rational constant used to prove
nonvanishing must be nonzero for each new power. We prove the first
point by retaining the identities at the integer zeros, and certify
the second by exact arithmetic at the fixed base index $96$.
The prime-density and phase bounds come from \cite{fifteen}. Their certificates
are in the \texttt{fifteen-beta-values/checks} directory of the project
repository~\cite{beta4repo}.

The underlying half-integer-pole construction and its local
valuation estimates come from Krisman \cite[Sections~1.2 and~4]{krisman}.
Differentiation to remove lower values follows the strategy discussed
in \cite{rivoal}; our direct source for its powered form is
\cite{powered}. The auxiliary-prime nonvanishing argument follows
\cite[Section~5]{fifteen} and \cite[Section~4]{powered}, which adapt
the prime-dilation strategy of Lai and Sprang
\cite[Lemma~2.3]{laisprang}.
The intended prerequisites are partial fractions, elementary
$p$-adic valuations, and asymptotic estimates for factorials.

\section{Why powers move the interval}\label{sec:forms}

The integer zeros of the kernel will do two jobs: they suppress the
early terms of the analytic sum and allow cancellation in its rational
constant. Its half-integer poles produce beta values. Increasing the
order of differentiation raises the first possible beta index, while
the pole orders determine the last.

Fix
\[
 N=96,\qquad r=\frac34,\qquad \delta_j=\frac{j}{96}
 \quad(0\le j\le30),
\]
and let $n$ range over positive multiples of $N$. Set
\[
 I_j(n)=[\delta_j n,(1-\delta_j)n]\cap\Z,\qquad
 B_a(n)=\#\{j:a\in I_j(n)\}\quad(0\le a\le n).
\]
We usually write $B_a$ for $B_a(n)$. At $n=96$, the blocks are
$[0,96],[1,95],\ldots,[30,66]$, and
$B_a(96)=1+\min(a,96-a,30)$. Thus the pole orders rise from $1$
at either end to $31$ in the middle. Define
\begin{equation}\label{eq:kernel}
 \Lambda_n=[-(1+r)n-1,rn]\cap\Z,\qquad
 R_n(t)=
 \frac{(2t+n+1)\prod_{\ell\in\Lambda_n}(t-\ell)^3}
 {\prod_{a=0}^n(t+a+\tfrac12)^{B_a}}.
\end{equation}
The two degree parameters we need are
\begin{equation}\label{eq:G}
 \gamma=\sum_{j=0}^{30}(1-2\delta_j)=\frac{341}{16},
 \qquad G=\gamma-3(1+2r)=\frac{221}{16}.
\end{equation}
There are $\gamma n+31$ denominator factors and
$3\#\Lambda_n=15n/2+6$ integer-zero factors. Consequently
$\deg R_n=-Gn-24$. The linear prefactor cancels one order at the
central pole; keeping the nominal orders $B_a$ simply gives a zero
leading partial-fraction coefficient there.

The kernel satisfies
\begin{equation}\label{eq:symmetry}
 R_n(-n-1-t)=R_n(t).
\end{equation}
Indeed, reflection permutes the integer zeros and the pole blocks.
The numerator product has even degree, whereas the prefactor and
the denominator both change sign.

Fix a positive integer $q$, and put
\begin{equation}\label{eq:L}
 L_{q,n}=\sum_{t=0}^{\infty}(-1)^t
                  \left(\frac{d}{dt}\right)^{2q}R_n(t)^q.
\end{equation}
Each integer zero now has order $3q>2q$. At integer arguments,
\begin{equation}\label{eq:integerzeros}
 (R_n^q)^{(2q)}(t)=0\quad(0\le t\le rn),\qquad
 (R_n^q)^{(2q)}(-\ell)=0\quad(1\le\ell\le n)
\end{equation}

Write
\[
 R_n(t)^q=\sum_{a=0}^n\sum_{i=1}^{qB_a}
               \frac{a_{q,n;a,i}}{(t+a+\tfrac12)^i},
 \qquad
 \sigma_a(s)=\sum_{h=0}^{a-1}\frac{(-1)^h}{(2h+1)^s}.
\]
Absent coefficients are set to zero, and
$(i)_d=i(i+1)\cdots(i+d-1)$ denotes the rising factorial.
Termwise differentiation and a shift of the beta series give
\begin{align}
 L_{q,n}&=\rho_{q,n}+\sum_{i=1}^{31q}\tau_{q,n;i}\beta(i+2q),
                                                       \label{eq:betaform}\\
 \tau_{q,n;i}&=2^{i+2q}(i)_{2q}
                         \sum_{a=0}^n(-1)^a a_{q,n;a,i},
                                                       \label{eq:tau}\\
 \rho_{q,n}&=-\sum_{a,i}(-1)^a2^{i+2q}(i)_{2q}
                         a_{q,n;a,i}\sigma_a(i+2q).     \label{eq:rho}
\end{align}
All series converge absolutely. Reflection
\eqref{eq:symmetry} gives
$a_{q,n;n-a,i}=(-1)^i a_{q,n;a,i}$.
Since $n$ is even, the terms with odd $i$ cancel in
\eqref{eq:tau}. Thus
\begin{equation}\label{eq:range}
 L_{q,n}\in\Q+\sum_{\substack{2q+2\le s\le 33q\\s\ {\rm even}}}
                                     \Q\beta(s).
\end{equation}
For $q=2$ the upper endpoint is $66$; for $q=3$ the nominal
endpoint $99$ is odd and disappears, leaving $98$.

\section{The same clearing factor works after powering}
\label{sec:arithmetic}

The denominator estimate in \cite{fifteen} discards an isolated
pole whenever it is the only source of a bad evaluation denominator.
We show that powering preserves this choice. This is essential:
a larger denominator estimate would need a new comparison with
analytic decay.

\subsection{The inherited integer}

We specify the integer $D_n$ of \cite[Section~3 and Appendix~A]{fifteen},
so that the arithmetic and analytic bounds refer to the same
normalization. Write $v_p(p)=1$ and $v_p(0)=+\infty$.
For an odd prime set
\[
 \lambda_p=\lfloor\log_p(7n)\rfloor,\qquad
 \nu_p(a)=3\sum_{\ell\in\Lambda_n}v_p(2a+1+2\ell)
                   -\sum_{b\ne a}B_bv_p(b-a).
\]
At $2$, put $Z=3\#\Lambda_n$ and
\[
 \lambda_2=\lfloor\log_2 n\rfloor,\qquad
 \nu_2(a)=1-Z-\sum_{b\ne a}B_bv_2(b-a).
\]
The small-prime exponents are
\begin{align}
 e_p&=\max\{0,\max_a(-\nu_p(a)+33\lambda_p)\}
                     &&(p\text{ odd},\ p^2\le7n),       \label{eq:esmall}\\
 e_2&=\max\{0,\max_a(-\nu_2(a)+31\lambda_2)\}.            \label{eq:e2}
\end{align}

For an odd prime $p\le n$ with $p^2>7n$, define
\begin{align*}
 N_{p,j}(a)&=\#\{b\in I_j(n):b\equiv a\pmod p\},\\
 c_p(a)&=3\#\{\ell\in\Lambda_n:p\mid2a+1+2\ell\}
                                  +\ind_{p\mid n-2a},\\
 F_p(a)&=\sum_{j=0}^{30}N_{p,j}(a)-c_p(a).
\end{align*}
A pole-address class is \emph{repeated} if it contains at least two
addresses in $\{0,\ldots,n\}$, and \emph{isolated} otherwise.
Let $M_p$ be the maximum of $F_p(a)$ over repeated classes, and put
\begin{equation}\label{eq:elarge}
 e_p=\max(0,M_p+2),\qquad
 D_n=\prod_{p\le n}p^{e_p}.
\end{equation}
For $p>n$ the exponent is zero.

\begin{proposition}\label{prop:clear}
For each fixed positive integer $q$, all the coefficients
$D_n^q\rho_{q,n}$ and $D_n^q\tau_{q,n;i}$ are integers.
\end{proposition}

\subsection{Local series and the small primes}

A convenient way to keep track of powering is the following elementary
rule. Say that a series has bound $(v,\lambda)$ at $p$ if its
degree-$h$ coefficient has valuation at least $v-\lambda h$.
Products add the first parameters and preserve a common slope
$\lambda$. Hence its $q$th power has bound $(qv,\lambda)$:
every summand in a coefficient of total degree $h$ has valuation
at least $qv-\lambda h$.

At $t=-a-\tfrac12+u$, remove the nominal pole of $R_n$. The remaining
series is
\begin{equation}\label{eq:local}
 H_a(u)=(n-2a+2u)
 \prod_{\ell\in\Lambda_n}(-a-\tfrac12-\ell+u)^3
 \prod_{b\ne a}(b-a+u)^{-B_b}.
\end{equation}
The coefficient $a_{q,n;a,i}$ is
$[u^{qB_a-i}]H_a(u)^q$.
The local estimates of \cite[Appendix~A]{fifteen}, obtained by
expanding the reciprocal linear factors, give the bounds
$(\nu_p(a),\lambda_p)$ and $(\nu_2(a),\lambda_2)$ for $H_a$.
Therefore
\begin{equation}\label{eq:poweredlocal}
 v_p(a_{q,n;a,i})\ge q\nu_p(a)-\lambda_p(qB_a-i).
\end{equation}
At an odd prime, a finite harmonic sum of exponent $i+2q$
costs at most $(i+2q)\lambda_p$. The combined lower bound is
\[
 q\nu_p(a)-\lambda_p(qB_a-i)-(i+2q)\lambda_p
 =q\{\nu_p(a)-(B_a+2)\lambda_p\}\ge-qe_p.
\]
At $2$, the harmonic denominators are odd, so no such cost occurs;
$B_a\le31$ and \eqref{eq:e2} suffice. The powers of two and the
integer rising factorial in \eqref{eq:tau}--\eqref{eq:rho}
only improve these bounds.

\subsection{Why an isolated pole remains harmless}

For odd $p$ with $p^2>7n$, the local calculation in
\cite[Section~3, local-coefficient lemma]{fifteen} gives $H_a$
bound $(B_a-F_p(a),1)$. Indeed, each nonunit linear factor has
valuation one, every nonzero difference is smaller than $p^2$,
and the central factor $2u$ also has bound $(1,1)$.
Consequently
\begin{equation}\label{eq:repeated}
 v_p(a_{q,n;a,i})\ge i-qF_p(a).
\end{equation}
If $a$ is isolated, every denominator factor in \eqref{eq:local}
has unit constant term, and the whole local series is integral at
$p$. This last assertion holds for every odd prime, regardless
of its size.

It remains to handle the rational correction. For $1\le\ell\le n$ set
\[
 T_q(\ell,a)=\sum_i(i)_{2q}a_{q,n;a,i}
                       (a-\ell+\tfrac12)^{-i-2q}.
\]
Reindexing the finite harmonic sums in \eqref{eq:rho}, and then
using \eqref{eq:integerzeros}, gives
\begin{equation}\label{eq:complement}
 \rho_{q,n}=-\sum_{\ell=1}^n(-1)^\ell
                      \sum_{a\ge\ell}T_q(\ell,a),
 \qquad
 \sum_{a\ge\ell}T_q(\ell,a)
             =-\sum_{a<\ell}T_q(\ell,a).
\end{equation}
The second identity lets us choose either side of $-\ell$.
For a fixed $\ell$, an evaluation denominator is divisible by
$p$ only in the single class $2(a-\ell)+1\equiv0\pmod p$.
If this class is isolated, choose the side of
\eqref{eq:complement} that does not contain its address.

For example, among addresses $0,\ldots,12$ modulo $11$ the class
of $8$ is isolated. At $\ell=3$ it is the only class with a bad
evaluation denominator. The sum over $a<3$ avoids it. The identity
making this possible still holds after powering because the
integer zero has order $3q$, greater than the derivative order $2q$.

If the bad class is repeated and $p^2>7n$, its evaluation
denominators have valuation one. Equation~\eqref{eq:repeated}
then gives
\[
 v_p\bigl((i)_{2q}a_{q,n;a,i}
                    (a-\ell+\tfrac12)^{-i-2q}\bigr)
 \ge i-qF_p(a)-(i+2q)
 \ge-q(M_p+2)\ge-qe_p.
\]
Terms with unit evaluation denominators, and the beta coefficients,
are covered by the same exponent; isolated coefficients are integral.
For $p>n$ all classes are isolated, so the complementary choice avoids
the bad class entirely, even if its denominators have higher
valuation. Together with the small-prime argument this proves
Proposition~\ref{prop:clear}.

\section{Decay with this normalization}\label{sec:decay}

The importance of Proposition~\ref{prop:clear} is that its integer is
the $q$th power of the original one. Both the main arithmetic cost
and the analytic decay are therefore multiplied by $q$.

In $\log D_n-G\log(n!)$, the large-prime contribution is at most
$(8.10809+o(1))n$ by the certificate \cite[(4.5)]{fifteen} for
the isolated-pole exponents \eqref{eq:elarge}. The small odd primes
contribute $o(n)$ and the prime $2$ contributes at most $15n\log2+o(n)$
\cite[Section~4 and Appendix~A]{fifteen}. Hence
\begin{equation}\label{eq:Dsize}
 \log D_n\le G\log(n!)
                +\left(8.10809+15\log2+o(1)\right)n.
\end{equation}
Put $H(u)=u\log u$ for $u>0$. For $x>r$, define
\[
 \Phi(x)=3\{H(x+1+r)-H(x-r)\}
          -\sum_{j=0}^{30}\{H(x+1-\delta_j)-H(x+\delta_j)\}.
\]
The uniform Stirling estimate of \cite[Section~4]{fifteen} gives,
for $t\ge rn+1$,
\begin{equation}\label{eq:stirling}
 \log R_n(t)=-G\log(n!)+n\Phi(t/n)+O(\log(n+t)).
\end{equation}
The global maximum of $\Phi$ is bounded in \cite[(4.9)]{fifteen},
and the derivative estimate in the same proof controls the tail. Thus
\begin{equation}\label{eq:phase}
 \Phi(x)\le-18.83511\quad(x>r),\qquad
 \Phi(x)\le\Phi(3)-3\log(x/3)\quad(x\ge3).
\end{equation}

\begin{proposition}\label{prop:decay}
For each fixed positive integer $q$,
\[
 \limsup_{\substack{n\to\infty\\96\mid n}}
       \frac1n\log|D_n^qL_{q,n}|
 \le q\bigl(8.10809+15\log2-18.83511\bigr)<-0.3298q,
\]
where $\log0=-\infty$.
\end{proposition}
\begin{proof}
For $t\ge rn+1$ every linear factor is bounded away from zero
and there are $O(n)$ factors, counted with multiplicity.
Each logarithmic derivative of fixed positive order of $R_n$
is therefore $O_q(n)$ on this range. The formula for a derivative
of an exponential gives
\[
 |(R_n^q)^{(2q)}(t)|\le C_q n^{2q}R_n(t)^q.
\]
The first identity in \eqref{eq:integerzeros} removes all preceding
integer terms.

For $rn+1\le t\le3n$, equations
\eqref{eq:stirling}--\eqref{eq:phase} bound the sum by a polynomial
in $n$ times $(n!)^{-qG}\exp(-18.83511qn)$.
For $t>3n$, the second bound in \eqref{eq:phase} supplies
$(3n/t)^{3qn}$. Its sum absorbs the fixed polynomial in $n+t$
coming from the Stirling error and again costs only a polynomial
in $n$. Combining with \eqref{eq:Dsize} cancels the factorial
terms and proves the assertion. The strict numerical inequality
follows from $\log2<0.6931472$, since
$\sum_{j=0}^9(0.6931472)^j/j!>2$.
\end{proof}

\section{A seed that prevents vanishing}\label{sec:seed}

Small integral forms give a contradiction only when they are
nonzero. We use a different increasing prime for each form.
After dilation, all beta coefficients will vanish modulo that prime,
while a fixed rational seed constant survives. This permits us
to avoid the finitely many primes dividing any hypothetical common
denominator of the beta values.

\subsection{Centering and the trimmed seed}

The change of variable $x=2t+n+1$ makes the poles integers.
Put
\[
 F_n(x)=
 \frac{x\prod_{\substack{v\ {\rm odd}\\|v|\le(5/2)n+1}}(x-v)^3}
      {\prod_{a=0}^n(x-n+2a)^{B_a}},
 \qquad
 F_n(x)^q=\sum_{a,i}\frac{c_{q,n;a,i}}{(x-n+2a)^i}.
\]
Counting powers of two gives
\begin{equation}\label{eq:center}
 R_n(t)=2^{Gn+25}F_n(2t+n+1),\qquad
 L_{q,n}=2^{q(Gn+25)+2q}\ell_{q,n},
\end{equation}
where the centered form has expansion
\begin{equation}\label{eq:centerform}
 \ell_{q,n}=\sum_{a,i}(-1)^a(i)_{2q}c_{q,n;a,i}
                         \{\beta(i+2q)-\sigma_a(i+2q)\}.
\end{equation}
Denote its rational constant by $\kappa_{q,n}$.

Take $n=Np=96p$, where $p>337$ is prime. The zeros divisible by
$p$ in $F_{Np}(pu)$ are $pv$ with $v$ odd and
$|pv|\le240p+1$. Thus $|v|\le239$.
The selected pole addresses are $pa$, $0\le a\le96$, and
$B_{pa}(Np)=B_a(96)$. They give the fixed rational function
\begin{equation}\label{eq:seedfunction}
 f(u)=
 \frac{u\prod_{\substack{v\ {\rm odd}\\|v|\le239}}(u-v)^3}
      {\prod_{a=0}^{96}(u-96+2a)^{B_a(96)}}.
\end{equation}
The endpoint $239$ is important: this is not $F_{96}$, whose
last positive zero is $241$.

For clarity, the factor counts are
\[
 W=\gamma N=2046,\qquad GN=1326,\qquad
 A_0=(W+31)-1-3\cdot240=1356.
\]
The last count says that $p^{A_0}F_{Np}(pu)$ has a nonzero
reduction after the selected factors have been extracted.

\begin{lemma}\label{lem:dilate}
Write
\[
 f(u)^q=\sum_{a=0}^{96}\sum_{i=1}^{qB_a(96)}
             \frac{\widetilde c_{q;a,i}}{(u-96+2a)^i}.
\]
For every prime $p>337$,
\begin{align}
 p^{1356q-i}c_{q,Np;pa,i}
       &\equiv(-1)^q\widetilde c_{q;a,i}\pmod p,
                                                     \label{eq:selected}\\
 v_p(c_{q,Np;a,i})&\ge i-1326q\qquad(p\nmid a).
                                                     \label{eq:unselected}
\end{align}
The quantities reduced in \eqref{eq:selected} are integral at $p$.
\end{lemma}
\begin{proof}
We include the counts from \cite[Section~5]{fifteen} to explain
why they pass to powers. In each block $[pj,p(96-j)]$, every
nonzero address class occurs $96-2j$ times and the zero class
occurs once more. Hence the weighted pole counts are $W$ for
each nonzero class and $W+31$ for zero.

After extracting the selected factors from $F_{Np}(pu)$, all
remaining factors reduce to constants. Their combined unit is $-1$.
Indeed, the $240p$ odd numbers from $-240p+1$ to $240p-1$
contain each residue $240$ times. Their nonzero residue product
is $1$ by Wilson's theorem. The two extra endpoints
$\pm(240p+1)$ contribute $-1$, also after cubing.
In the denominator, each complete nonzero residue block contributes
$-1$, while powers of $2^{p-1}$ are $1$; the total sign is
$(-1)^W=1$. Thus, locally away from the other selected poles,
\begin{equation}\label{eq:functionreduction}
 p^{1356}F_{Np}(pu)\equiv-f(u)\pmod p.
\end{equation}
At the pole $u=96-2a$, multiply by its nominal pole-removing
power. The other selected differences $2(b-a)$ are units, as
are the constant terms of unselected denominator factors after
their factors of $p$ have been separated. The resulting local
series is in $\Z_{(p)}[[w]]$.
Raising it to the $q$th power and extracting the coefficient
of degree $qB_a(96)-i$ proves \eqref{eq:selected}.

At a nonselected address the prefactor is a unit, the numerator
has at least $240$ zeros in each residue class before cubing,
and its pole-address class has weighted count $W$. Removing
the nominal pole gives a local bound
$(B_a+720-W,1)=(B_a-1326,1)$.
All nonzero differences have valuation at most one: their
absolute values are below $337p<p^2$.
The power rule and extraction of degree $qB_a-i$ now give
\eqref{eq:unselected}.
\end{proof}

\subsection{Separating the constant from the beta coefficients}

Define the rational seed constant
\begin{equation}\label{eq:kq}
 k_q=-\sum_{a=0}^{96}\sum_i(-1)^a(i)_{2q}
                    \widetilde c_{q;a,i}\sigma_a(i+2q),
 \qquad \tau_q=(1356+2)q=1358q.
\end{equation}
The two parts of $\tau_q$ match the pole and the derivative. At a
selected address $pa$, multiplying the term of order $i$ by $p^{\tau_q}$
gives $p^{i+2q}$ times the $p$-integral quantity on the left of
\eqref{eq:selected}. A harmonic denominator of exponent $i+2q$ can
consume exactly that power. At an unselected address,
\eqref{eq:unselected} leaves at least $(1356-1326)q=30q$ further
powers of $p$ after the same loss. This is why only the selected
addresses can contribute to the seed.

\begin{proposition}\label{prop:unit}
Every beta coefficient of $p^{\tau_q}\ell_{q,Np}$ is divisible
by $p$, and all its coefficients are integral at $p$. Its rational
constant satisfies
\begin{equation}\label{eq:constantreduction}
 p^{\tau_q}\kappa_{q,Np}
             \equiv(-1)^q\chi_4(p)k_q\pmod p,
 \qquad \chi_4(p)=(-1)^{(p-1)/2}.
\end{equation}
\end{proposition}
\begin{proof}
Put $s=i+2q$. The scaled beta coefficients have valuation at least
$s$ by the preceding bounds and the integrality of $(i)_{2q}$.
Every $2h+1$ with $h<Np$ is below $2Np<p^2$, so a harmonic term
costs at most $s$ powers of $p$. This proves integrality of the
scaled constant and makes its unselected terms zero modulo $p$.
For a selected address $pa$, only denominators divisible by $p$
survive in $p^s\sigma_{pa}(s)$. Writing
$2h+1=p(2v+1)$ gives
\[
 h=pv+\frac{p-1}{2},\quad 0\le v<a,\qquad
 p^s\sigma_{pa}(s)\equiv\chi_4(p)\sigma_a(s)\pmod p.
\]
There is no further $p$ in $2v+1$, since $p>337>191$.
Combining this identity, $(-1)^{pa}=(-1)^a$, and
\eqref{eq:selected} with the constant part of
\eqref{eq:centerform} proves \eqref{eq:constantreduction}.
\end{proof}

\section{Exact finite certificates and the conclusion}
\label{sec:certificate}

It remains to prove $k_2\ne0$ and $k_3\ne0$. This is a calculation
on the fixed function \eqref{eq:seedfunction}, not on an increasing
sequence of kernels. At $u_a=96-2a$, form
\[
 J_a(w)=(u_a+w)
 \prod_{\substack{v\ {\rm odd}\\|v|\le239}}(u_a-v+w)^3
 \prod_{\substack{0\le b\le96\\b\ne a}}
                       (2(b-a)+w)^{-B_b(96)}.
\]
Then
\[
 \widetilde c_{q;a,i}=[w^{qB_a(96)-i}]J_a(w)^q.
\]
The exact reductions of \eqref{eq:kq} are
\begin{center}
\begin{tabular}{@{}ccccc@{}}
\toprule
$q$ & $k_q\bmod10007$ & $k_q\bmod10009$ & $k_q\bmod10037$
    & Coefficients per prime\\
\midrule
$1$ & $3738$ & $5333$ & $938$ & $2077$\\
$2$ & $8279$ & $5481$ & $5660$ & $4154$\\
$3$ & $9978$ & $6289$ & $4500$ & $6231$\\
\bottomrule
\end{tabular}
\end{center}
The first two residues at $q=1$ reproduce \cite[Section~5]{fifteen}.
The column for $10037$ comes from an independent reconstruction.
Each of the last two rows contains a nonzero residue, proving
$k_2\ne0$ and $k_3\ne0$, respectively.

Here is a reproducible specification of the calculation.
One route multiplies the truncated factors in $J_a^q$ directly.
The other writes its unit part as $g(w)=\prod_\nu(v_\nu+w)^{e_\nu}$,
with integer exponents, and uses
\[
 \frac{g'(w)}{g(w)}
  =\sum_{j\ge0}\left(\sum_\nu
                  e_\nu\frac{(-1)^j}{v_\nu^{j+1}}\right)w^j.
\]
If the coefficient in parentheses is $d_j$ and
$g(w)=\sum_{k\ge0}g_kw^k$, then
\[
 g_0=\prod_\nu v_\nu^{e_\nu},\qquad
 g_k=\frac1k\sum_{j=0}^{k-1}d_jg_{k-1-j}.
\]
At the central pole, remove the factor $w^q$ before this
recurrence and restore it afterwards.
All three moduli are prime by trial division. Pole differences
are at most $192$, nonzero numerator differences at most $335$,
harmonic denominators at most $191$, and the recurrence indices
are below $93$. All inversions used above are therefore legitimate.

The two methods agree on every coefficient. The seed checker also verifies
reflection and the vanishing odd beta coefficients and computes the constant
in two summation orders. From the \texttt{beta-intervals/} directory of the project
repository~\cite{beta4repo}, run
\begin{verbatim}
python3 -B checks/interval_seeds.py
python3 -B checks/interval_seeds.py --cross-check
\end{verbatim}
The second command requires \texttt{python-flint}; the first uses
only the Python standard library.

\begin{proof}[Proof of Theorem~\ref{thm:main}]
Fix $q=2$ or $q=3$, and suppose that all beta values in
\eqref{eq:range} are rational. One positive integer $T$ clears
their denominators. Proposition~\ref{prop:clear} gives
\[
 T D_n^qL_{q,n}\in\Z\qquad(96\mid n),
\]
and Proposition~\ref{prop:decay} makes these integers tend to zero.

To show that infinitely many of these integers are nonzero, we use
$p^{\tau_q}$ at the single prime $p$ on the subsequence $n=96p$.
The corresponding seed $k_q$ is a nonzero rational number.
For every sufficiently large prime outside the finite set dividing
$T$ or the numerator or denominator of $k_q$, all the beta values
are integral at $p$. Proposition~\ref{prop:unit} makes their
coefficients in $p^{\tau_q}\ell_{q,96p}$ divisible by $p$ and
its constant a unit. Thus $\ell_{q,96p}\ne0$, and
\eqref{eq:center} gives $L_{q,96p}\ne0$. Infinitely many
nonzero integers $T D_{96p}^qL_{q,96p}$ would tend to zero,
a contradiction.
\end{proof}

The argument up to the finite calculation applies to every fixed
$q$ for which $k_q\ne0$. The two assertions above use only the
separately certified cases $q=2,3$.

\paragraph{Verification.}
A second AI model independently audited the original proof and exposition of
Theorem~\ref{thm:main} and reproduced the seeds and decay bounds.

\begin{thebibliography}{9}
\bibitem{fifteen}
twinheterodox,
\emph{At least one of $\beta(4),\beta(6),\ldots,\beta(32)$ is irrational},
2026, \href{https://hexagonmath.org/2611.00346v1}{Hexagon 2611.00346v1}.

\bibitem{powered}
twinheterodox,
\emph{Irrational even beta values from powered hypergeometric kernels},
2026, \href{https://hexagonmath.org/2611.00336v1}{Hexagon 2611.00336v1}.

\bibitem{krisman}
S.~Krisman,
\emph{One of $\beta(4),\beta(6),\ldots,\beta(36)$ is irrational},
2026, \href{https://hexagonmath.org/2610.00285v2}{Hexagon 2610.00285v2}.

\bibitem{laisprang}
L.~Lai and J.~Sprang,
\emph{Many $p$-adic odd zeta values are irrational},
2025, \href{https://arxiv.org/abs/2306.10393v2}{arXiv:2306.10393v2}.

\bibitem{rivoal}
T.~Rivoal,
\emph{Irrationalit\'e d'au moins un des neuf nombres
$\zeta(5),\zeta(7),\ldots,\zeta(21)$},
Acta Arith. \textbf{103} (2002), 157--167,
\href{https://doi.org/10.4064/aa103-2-5}{doi:10.4064/aa103-2-5}.

\bibitem{beta4repo}
twinheterodox, \emph{Irrationality and linear independence of Dirichlet beta values},
GitHub repository, 2026. \url{https://github.com/twinheterodox/beta4}.
\end{thebibliography}
\end{document}
